\section*{Bab \ref{ch:g2:mult} --- Perkalian: Langkah Pertama}

\begin{solution}{exo:g2:mult:1}
$2 + 2 + 2 = 3 \times 2 = 6$; \quad
$5 + 5 + 5 + 5 = 4 \times 5 = 20$; \quad
$10 + 10 + 10 = 3 \times 10 = 30$.
\end{solution}

\begin{solution}{exo:g2:mult:2}
$3$ baris berisi $5$: $3 \times 5 = 15$. Kalau dimiringkan, $5$ baris
berisi $3$: $5 \times 3 = 15$ --- $15$ yang sama.
\end{solution}

\begin{solution}{exo:g2:mult:3}
Tabel $2$: $2, 4, 6, 8, 10, 12$. Tabel $5$:
$5, 10, 15, 20, 25, 30$.
\end{solution}

\begin{solution}{exo:g2:mult:4}
$2 \times 7 = 14$; \quad $5 \times 3 = 15$; \quad
$10 \times 4 = 40$; \quad $2 \times 9 = 18$; \quad
$5 \times 6 = 30$.
\end{solution}

\begin{solution}{exo:g2:mult:5}
Kembar: $12$; $16$; $20$. Separuh: separuh dari $12$ adalah $6$;
separuh dari $20$ adalah $10$.
\end{solution}

\begin{solution}{exo:g2:mult:6}
$3, 6, 9, 12, 15, 18$. Jadi $3 \times 5 = 15$.
\end{solution}

\begin{solution}{exo:g2:mult:7}
$4 \times 5 = 20$ kursi.
\end{solution}

\begin{solution}{exo:g2:mult:8}
$3 \times 8 = 24$ kaki ($8, 16, 24$).
\end{solution}

\begin{solution}{exo:g2:mult:9}
$2$ baris berisi $6$: $2 \times 6 = 12$. $3$ baris berisi $4$:
$3 \times 4 = 12$.
\end{solution}

\begin{solution}{exo:g2:mult:10}
Bilang loncat $5$: $5, 10, 15, 20$ --- tiga kotak hanya memberi $15$
(terlalu sedikit), empat kotak memberi $20$. Mia harus membeli $4$ kotak.

\end{solution}
